\section{Forty-one contacts with fixed coefficients}\label{sec:sp-contacts}

Set
\[
 \mathcal D=\{0,1,2401,-343/32\},\qquad
 R=47027=31\cdot37\cdot41,
\]
\[
 W=25878772920,\qquad c=18484837801,
 \qquad V_0=32W=828120733440.
\]
For $k\le9$ rational targets $a_1,\ldots,a_k$, put
$A_i=Ra_i+3$. With common variables $u,v$, set $w_i=A_iv-u$.
The contact list is
\begin{equation}\label{eq:sp-fixed-contacts}
 \frac v{V_0},\qquad u-\delta v\ (\delta\in\mathcal D),
 \qquad -(w_i-\delta v)\
       (\delta\in\mathcal D,\ 1\le i\le k).
\end{equation}
It contains $5+4k$ entries. Every entry is required to be nonzero.
Its normalizing denominators have prime support in $S_*$. Here
$M(u:v)=M(u/v)$: at each prime $q$ the five contact-form valuations
are normalized by subtracting $\min(v_q(u),v_q(v))$. The radical
records odd normalized valuations outside $S_*$.

\begin{lemma}\label{lem:sp-contacts}
For every list of at most nine integers $a_i$, there are integers $u,v$
for which all entries of \eqref{eq:sp-fixed-contacts} are positive
ordinary primes exceeding any prescribed bound and congruent to $c$
modulo $W$.
\end{lemma}
\begin{proof}
Use the lattice coset
\begin{equation}\label{eq:sp-lattice}
 v=32W(c+Wy),\qquad u=c+Wx,\qquad x,y\in\mathbb Z.
\end{equation}
Every divided form in \eqref{eq:sp-fixed-contacts} is integral and
congruent to $c$ modulo $W$. Indeed $v/32$ is divisible by $W$,
and $A_i$ is integral. The integer $c$ is $1$ modulo $8$ and
modulo every odd prime at most $29$ other than $7$, where it is $3$.
These are the required initial prime and successor conditions.

Duplicate targets may be removed when applying the counting theorem;
the resulting values also satisfy their repeated occurrences. For distinct
targets, all forms are pairwise nonproportional over $\mathbb Q$.
Apart from $v/V_0$, their slopes are
\[
 \mathcal S=\mathcal D\cup\bigcup_i(A_i-\mathcal D).
\]
The only nonzero integral differences of elements of $\mathcal D$ are
$\pm1,\pm2400,\pm2401$, none divisible by $31$. Thus distinct
secondary families cannot meet. The integral sums of two elements of
$\mathcal D$ are
\[
 0,1,2,2401,2402,4802.
\]
Their residues modulo $31$ are $0,1,2,14,15,28$, so none equals
$A_i\equiv3\pmod{31}$. Hence no secondary slope is a primary slope.

We verify local admissibility, including the successor exclusions.
Fix a prime $p>29$. The map in \eqref{eq:sp-lattice} is invertible
modulo $p$. Write $\gamma=V_0^{-1}\in\mathbb F_p$. For
$p\in\{31,37,41\}$, all $A_i$ are $3$ modulo $p$, so the
image of $\mathcal S$ has at most eight elements. For $p\ge43$
it has at most forty elements. Choose distinct
$b,e\in\mathbb F_p\setminus\mathcal S$. Reverse their order if
necessary to ensure $e-b\ne-\gamma$, and put
\[
 v=(e-b)^{-1},\qquad u=bv.
\]
For every $s\in\mathcal S$,
\[
 u-sv=v(b-s)\ne0,\qquad u-sv+1=v(e-s)\ne0.
\]
Also $\gamma v$ and $\gamma v+1$ are nonzero. The reversal is
possible because $p$ is odd and $\gamma\ne0$. This proves the
required simultaneous nonvanishing at every prime above $29$.

There is a nonempty real cone on which all the forms are positive:
take $v>0$ and $u/v>\max_i|A_i|+2413$. Choose a closed box in
this cone. The source prime-pattern theorem applies to the fixed affine
system on this box and the lattice coset. To use the source fiber sieve,
add an unused third coordinate. A fiber of one contact form is then a
two-dimensional affine plane. The other zero loci on it are distinct
lines, outside the fixed determinant primes and the prime value defining
the fiber. This follows because the original forms are homogeneous
and pairwise nonproportional in $(u,v)$. The fixed determinant primes
are included in the initial cutoff. Thus the source uniform fiber
estimate applies, with constants depending on the fixed targets.

More explicitly, if $m$ is the number of distinct forms and $L$ is
the initial sieve cutoff, the source lower estimate has the form
\[
 \frac{c_1 X^3}{(\log X)^m(\log L)^{c_2}},
\]
whereas the patterns failing the ordinary-prime order conditions have
upper bound
\[
 O\!\left(\frac{X^3}{L(\log X)^m}\right)
       +o\!\left(\frac{X^3}{(\log X)^m}\right).
\]
The constants are fixed after the target list and box are fixed.
The finitely many small-prime choices above preserve positivity of the
singular product; for large primes its unmodified factors are
$1+O(p^{-2})$, and the moving successor exclusions give the displayed
fixed power of $\log L$. Choose $L$ sufficiently large, then $X$
sufficiently large. A surviving pattern has all the required ordinary
prime values. The box may be dilated so that those values exceed any
prescribed bound.
\end{proof}

\begin{lemma}\label{lem:sp-shared-height}
Suppose every contact radical belonging to the slopes $u/v,w_i/v$
in \eqref{eq:sp-fixed-contacts} has its $K$th power dividing a nonzero
integer $N$. Then
\[
 h(a_i)<2H|N|^{24/K}+14.
\]
If the targets include a common nonzero denominator $e$ and the two
coefficients of $e\,x(Q)$ for a point $Q\in G$, then
\begin{equation}\label{eq:sp-shared-height}
 h(x(Q))<7H|N|^{24/K}.
\end{equation}
\end{lemma}
\begin{proof}
Since $A_i=u/v+w_i/v$, the contact-height bound gives
\[
 h(a_i)\le2H|N|^{24/K}+2\log2+\log3+\log R
         <2H|N|^{24/K}+14.
\]
Writing $e\,x(Q)=a+b\sqrt2$, one has
\[
 h(x(Q))\le h(a)+h(b)+h(e)+h(\sqrt2)+\log2.
\]
The additive term after substitution is less than $50$, and the fixed
constant $H$ is greater than $50$. This proves
\eqref{eq:sp-shared-height}.
\end{proof}

\section{One simultaneous comparison at precision sixty-four}
\label{sec:sp-comparison}

Fix nonzero $Q_i=n_iP_*$ for $1\le i\le4$, rational proxies $b_i$,
and a common nonzero $J=jP_*$. Put $Q_0=P_*$. At each contact choose
nonzero points $X_0=hP_*$ and $X_i=m_iP_*$, and set $Y_i=X_i-4Q_i$.
The unequal-abscissa subtraction charts are imposed. Use the truncated logarithm $L$ from~\eqref{eq:log}.
Test order one for the eleven field expressions
\[
 f(X_i)\ (0\le i\le4),\qquad f(J^\sigma),\qquad
 f(Y_i^\sigma)\ (0\le i\le4),
\]
and order $64$ for the nine expressions
\[
 \frac{L(z(X_i))}{L(z(X_0))}-b_i
       \ (1\le i\le4),\qquad
 \frac{z(Y_i^\sigma)}{z(J^\sigma)}\ (0\le i\le4).
\]
The displayed denominators have inverse witnesses. For a field test
$a+b\sqrt2$ of order $L$, use the rational tests
$(\theta^2-2)b$ and $a+\theta b$, both of order $L$, together
with $\theta^2-2$ at order one. Thus there are $41$ rational
tests. A fixed binary norm batch, including the constant leaf $1$,
imposes these tests at every prime outside $S_*$ where the contact
value has odd valuation.

\begin{lemma}\label{lem:sp-comparison}
At each such prime $q$, the comparison certificate implies
$b_i\in\mathbb Z_q$ and $b_i\equiv n_i\pmod{q^{64}}$.
If $b_i=n_i$, any finite collection of these certificates at ordinary
primes can be satisfied simultaneously.
\end{lemma}
\begin{proof}
The root test selects a place above $q$. The conjugate parameter ratios
give
\[
 h\equiv4\pmod{q^{64}},\qquad
 m_i\equiv4n_i\pmod{q^{64}},
\]
because $j$ is an integer. In particular $h$ is a $q$-unit.
The formal logarithm tail begins in degree $65$, and the source formal
expansion gives $v(\ell(z)-L(z))\ge65v(z)$ at these primes.
Because $h$ is a unit, $v(z(X_i))\ge v(z(X_0))$. Dividing the
tail estimate by the parameter of $X_0$ therefore gives
\[
 \frac{L(z(X_i))}{L(z(X_0))}
       \equiv m_i/h\pmod{q^{64}}.
\]
Here integrality of the index ratio follows from $h$ being a unit;
all truncated coefficient denominators are units since $q>76176$. The primary
tests now prove the assertion.

For completeness, let $h_q,j_q$ be the source ordinary local witnesses
at each selected place. Choose a positive common multiple $j$ of the
nonzero $j_q$. The congruences
\[
 t_q\equiv1\pmod{|j_q|},\qquad
 h_qt_q\equiv4\pmod{q^{64+v_q(j)}}
\]
are compatible, since the original local witnesses give
$h_q\equiv4\pmod{q^{v_q(j_q)}}$. Choose a solution with
$h_qt_q\notin\{0,4,-4\}$. Set
\[
 J=jP_*,\qquad X_0=h_qt_qP_*,\qquad X_i=h_qt_qn_iP_*.
\]
The first congruence preserves the required formal subgroups; the
second gives order $64$ relative to $J$. The excluded values ensure
every subtraction chart is valid, because $n_i\ne0$. Choose consistent
places at repeated primes and then choose the common $\theta$ by CRT.
The norm values are units at the corresponding positive prime contacts,
so the specialized parity completeness supplies the remaining witnesses.
\end{proof}

\section{Lucas growth and the defining relation}\label{sec:sp-lucas}

\begin{lemma}\label{lem:sp-lucas}
Suppose $s,t,z\in\mathbb Z$, $|s|=a\ge7$, $z\ne0$, and
\begin{equation}\label{eq:sp-lucas}
 t^2-s(s-2)tz+(s-2)^2z^2=1.
\end{equation}
Then $|t|>\exp(a-3)$. Conversely, for every such $s$ there are
positive integers $t,z$ satisfying \eqref{eq:sp-lucas}.
\end{lemma}
\begin{proof}
Put $w=\operatorname{sgn}(s)(s-2)z$. Then
$t^2-atw+w^2=1$. Let
$\rho=(a+\sqrt{a^2-4})/2$, and define
$U_0=0$, $U_1=1$, $U_{k+1}=aU_k-U_{k-1}$.
Every positive norm-one unit of $\mathbb Z[\rho]$ is a power of
$\rho$: multiply it by a power of $\rho$ to obtain
$v\in[1,\rho)$, write $v=A+B\rho$, and observe that
\[
 B=\frac{v-v^{-1}}{\rho-\rho^{-1}}\in[0,1)\cap\mathbb Z.
\]
Thus $B=0$ and $v=1$. Apply this to either $t-\rho w$ or its
negative. Since $w\ne0$, the identities
\[
 \rho^k=U_k\rho-U_{k-1},\qquad
 \rho^{-k}=U_{k+1}-U_k\rho
\]
give some $k\ge1$ with $|w|=U_k$ and
$|t|\in\{U_{k-1},U_{k+1}\}$.

For $s>0$, reduction modulo $a-2$ gives $U_k\equiv k$, so
$k\ge a-2$. For $s<0$, reduction modulo $a+2$ gives
$U_k\equiv(-1)^{k-1}k$, so $k\ge a+2$.
Induction gives $U_k\ge(a-1)^{k-1}$, hence
\[
 |t|\ge(a-1)^{a-4}>\exp(a-3).
\]
For the last inequality use $a-1\ge6>e^{3/2}$ and
$3(a-4)/2\ge a-3$. Conversely take $t=U_{k+1}$ and
$z=U_k/(a-2)$ with $k=a-2$ if $s>0$; take
$z=U_k/(a+2)$ with $k=a+2$ if $s<0$.
\end{proof}

For the rational input $X$, set $\beta=X^2+2$. Choose nonzero
$A=nP_*,S=sP_*,T=tP_*,Z=zP_*$. For the anchor, let $\mathcal T(t)$ be the seven-witness formula
\[
 u_0^2-2u_1^2-3u_2^2+6u_3^2=1,\qquad
 (t-2u_0)^2-8v_1^2-12v_2^2+24v_3^2=4.
\]
The source identifies it with $\mathbb Z_{(2)}\cap\mathbb Z_{(3)}$.
Set $C=3^{2^{24}}$, specified by $c_0=3$ and
$c_{i+1}=c_i^2$ for $0\le i<24$. If
$f(S)=f_0+f_1\sqrt2$, impose
\[
 \mathcal T(f_0/C),\qquad
 \mathcal T(f_1/C).
\]
The source anchor applies to $\lambda^2 f(S)$. Multiplication by
$\lambda^2$ acts on coefficient pairs by the integral matrix
$\left(\begin{smallmatrix}17&24\\12&17\end{smallmatrix}\right)$
of determinant one. Thus the displayed conditions are equivalent to
the original anchor conditions. The source local
calculation gives $|s|\ge3^{2^{23}-1}$, and the allowed indices
contain a nonzero subgroup of $\mathbb Z$. Thus the anchor may be
combined with any finite formal divisibilities. Its cutoff is
\begin{equation}\label{eq:sp-cutoff}
 |s|>(20H)^2,
\end{equation}
because
\[
 3(2^{23}-1)=25165821>4(5+4783227)=19132928
\]
and $\log_2 3>3/2$.

The separate scalar-height batch uses
\[
 r(\beta)=\beta(\beta-1)(\beta-2401)(32\beta+343)(2\beta-1)
\]
and the two coefficients of each of $f(J_0),f(S),z(S)/z(J_0)$,
in orders $1,1,1,1,48,48$. Its parity certificate and prescribed local
divisibilities give
\begin{equation}\label{eq:sp-scalar-height}
 h(\beta)\le H\sqrt{|s|}.
\end{equation}
Indeed every odd contact of $\beta$ outside $S_*$ gives an odd
valuation of $r(\beta)$; at a pole the valuation is multiplied by
five. The two coefficient tests make both points formal at every place
above that prime, and the parameter ratio implies $M(\beta)^{48}\mid s$.
The contact-height input gives \eqref{eq:sp-scalar-height}.
For integral $X$, $r(\beta)$ is a nonzero integer: $\beta\ge2$,
and $\beta=2401$ would require $X^2=2399$. The same exclusion holds
for rational $X$, since $48^2<2399<49^2$. The scalar batch is complete
by the stated finite local-divisibility property.

Choose a common nonzero denominator $e$ and the two rational coefficients
of each of $e\,x(A)$, $e\,x(S)$, $e\,x(T)$, and $e\,x(Z)$.
Apply the shared representation
\eqref{eq:sp-fixed-contacts} to these nine scalars. At each of its
forty-one contacts apply Lemma~\ref{lem:sp-comparison} to
$(A,S,T,Z)$ and proxies $(\beta,\xi,\tau,\zeta)$, and impose
\begin{equation}\label{eq:sp-proxy-lucas}
 \tau^2-\xi(\xi-2)\tau\zeta+(\xi-2)^2\zeta^2=1.
\end{equation}
All predicates here have the fixed polynomial encodings given elsewhere;
there are no quantified integers or primes.

\begin{theorem}\label{thm:sp-definition}
Under the arithmetic inputs in Section~\ref{sec:inputs}, the preceding
existential relation holds for $X\in\mathbb Q$ if and only if
$X\in\mathbb Z$.
\end{theorem}
\begin{proof}
Suppose first that the relation holds. Put
$B=\max(|n|,|s|,|t|,|z|)$ and
\[
 \Delta=t^2-s(s-2)tz+(s-2)^2z^2-1.
\]
The proxy congruences imply $M^{64}\mid\Delta$ for each coordinate
contact radical. If $\Delta\ne0$, then $B\ge|s|>7$ and
\[
 |\Delta|\le2B^4+6B^3+5B^2+1<5B^4.
\]
At an index attaining $B$, \eqref{eq:canonical} and
Lemma~\ref{lem:sp-shared-height} give
\[
 B^2-6<7H|\Delta|^{3/8}<14HB^{3/2}.
\]
Consequently $B<(14H+6)^2\le(20H)^2$, contrary to
\eqref{eq:sp-cutoff}. Thus \eqref{eq:sp-lucas} holds for the
integer indices, and Lemma~\ref{lem:sp-lucas} gives
$|t|>\exp(|s|-3)$.

If $\beta\ne n$, let $N=|\operatorname{num}(\beta-n)|$,
where the numerator is taken in reduced form. Local integrality and
the comparison congruences give $M^{64}\mid N$ at every coordinate
slot. Hence
\[
 B^2<7HN^{3/8}+6,\qquad
 N\le e^{h(\beta)}(1+|n|)\le2e^{h(\beta)}B.
\]
Since $7\cdot2^{3/8}<10$, it follows that
\[
 \log B<\frac8{13}\log(16H)+\frac3{13}h(\beta).
\]
Write $a=|s|>(20H)^2$. Then
$h(\beta)\le H\sqrt a<a/20$ and
$\log(16H)<16H<a/20$. Thus
$\log B<11a/260<a-3$, contradicting Lucas growth.
Therefore $\beta=n\in\mathbb Z$. For rational $X$, integrality
of $X^2$ implies integrality of $X$, by reducing its denominator.

Conversely let $X\in\mathbb Z$, so $\beta=X^2+2\ge2$, and set
$A=\beta P_*$. Choose the scalar-height witnesses, then an anchored
$S=sP_*$ satisfying their finite formal divisibilities. Choose positive
$t,z$ from Lemma~\ref{lem:sp-lucas} and put $T=tP_*,Z=zP_*$.
Choose an integer $e\ne0$ clearing all eight abscissa coefficients.
Lemma~\ref{lem:sp-contacts} supplies the shared ordinary prime pattern.
Set the proxies equal to $(\beta,s,t,z)$, and use
Lemma~\ref{lem:sp-comparison} for the local witnesses, choosing the
common root approximation last. Every remaining variable is an indicated
coordinate, inverse, or fixed arithmetic intermediate. All equations and
guards hold. No later choice changes an earlier target or anchor point.
\end{proof}
